% 6 Summary and conclusions (all). Verdicts as presented (deck slide 37): keep them unless new evidence changes them.
\section{Summary and conclusions}\label{sec:conclusions}

\subsection{What we expected, and what we found}
The plan stated three expected results. \Tab{tab:verdicts} sets each against the evidence. The first is only
\emph{partly confirmed}: concentration is now confirmed at country level (Gini 0.94), but the claim that capacity sits
far from its water source rests on two case studies, not a systematic measurement. The second, the efficiency
explanation, is \emph{partly supported}: the timing fits the end of building-efficiency gains, but chip efficiency was
not tested. The third is \emph{supported as a framework}: impacts do sort by scale, but several cells of the matrix are
documented rather than measured.

\begin{table}[t]\centering
\caption{Expected results (from the plan) against what we found.}
\label{tab:verdicts}
\footnotesize\renewcommand{\arraystretch}{1.15}
\begin{tabular}{@{}P{2.0cm}P{3.75cm}P{1.6cm}@{}}
\toprule
Expected & Found & Verdict \\
\midrule
Capacity is far more concentrated than population; part of it sits in water-stressed basins, often far from its
water source & 83\% of tracked AI power is in the US; country Gini 0.94; LQ about 20 in the US and Malaysia. 34\% of
capacity in high-stress basins today, 39--45\% of the 2030 plan. Distance to source shown for two cases only (231 and
193\,km in a straight line) & \textbf{Partly confirmed} \\
Cloud-era growth was absorbed by one-off efficiency gains, now used up, so AI-era growth outpaces efficiency &
Compute broke upward in 2011, yet data-centre electricity rose only 6\% in 2010--18. PUE flat since 2014; the post-2018
jump (7 to 18\% a year) matches AI-server demand ($<$2 to $>$40\,TWh). Chip efficiency not tested & \textbf{Partly
supported} \\
Electricity and water impacts are local or regional; carbon is global; minerals and e-waste are displaced & Water:
$<$0.1\% of the global limit, but 12 of 12 months over supply in Mesa's basin. Grid carbon varies 2.6-fold. Minerals
come from a few producers; e-waste collection is very uneven (0.7--42.8\%) & \textbf{Supported as a framework} \\
\bottomrule
\end{tabular}
\end{table}

\subsection{Limitations}\label{sec:limits}
The findings are only as good as the data behind them. In rough order of how much they matter:
\begin{enumerate}[leftmargin=1.3em,itemsep=1.5pt]
\item \textbf{The sample is AI-only and US-heavy.} Epoch AI tracks 96 large AI sites, 77 of them in the United States.
It is not a census of data centres: the IEA's 415\,TWh covers all of them. Shares, the Gini and the location quotients
describe where large AI campuses are being built, and countries such as India appear only through a case study.
\item \textbf{Water stress counts power-plant withdrawals.} Aqueduct's annual stress is total withdrawals over renewable
supply, and withdrawals include thermoelectric cooling water, most of which is returned. Basins with large power plants,
such as parts of Indiana and western New York, probably score ``high'' partly for that reason.
\item \textbf{Consumption against a withdrawal limit.} Our water figures are consumption; Earth3's freshwater boundary is
set on withdrawals. The comparison is conservative (withdrawals would be larger), but the two are not the same
quantity.
\item \textbf{US-derived grid water factors.} The indirect range uses 1.04--4.52\,L/kWh, from IEA- and LBNL-based values
that reflect the US grid. Applying them to India (Visakhapatnam) or Malaysia is an approximation; local generation mixes
could give lower or higher values.
\item \textbf{The load factor is assumed.} We assume sites draw 70--90\% of facility power on average. No site reports
this, and partly built campuses may run lower; water scales linearly with it.
\item \textbf{The seven-group Gini was a lower bound.} The deck's 0.927 pooled countries; our country-level 0.94 is the
exact value for the sample. It still cannot see inequality within countries.
\item \textbf{Breakpoints depend on the window.} World energy breaks in 1974 over 1950--2019 but in 2011 over
1990--2019, and 1950 cannot be tested because earlier data are decadal.
\item \textbf{Data-centre growth rests on published points.} Each data-centre row of \Tab{tab:growth} rests on two
published values, not a series we fitted, so it has no interval; the US rates are those stated by LBNL. Masanet et al.'s
+6\% is taken from the LBNL summary of the paper.
\item \textbf{The siting reading rests on public sources.} Section~\ref{sec:siting} uses statements, official documents and press reports,
some from campaigners and some from the company; Google's water assessment for Visakhapatnam is unpublished. Our air-cooling arithmetic uses
US-derived grid factors and a 10--20\% energy penalty taken from a company statement. We did not compare Visakhapatnam with other Indian hubs.
\item \textbf{Facility versus IT power.} Epoch's MW are facility power; we convert to IT energy with the PUE range. Sites
that report water do not state the basis, so the 0.001--7\,L/kWh spread is indicative only.
\end{enumerate}

\subsection{Responses matched to the scale of impact}\label{sec:responses}
Because each pressure has its own scale, no single policy covers them.
\textbf{Locally}, siting rules can keep large campuses out of basins in extreme stress or require non-potable water
(reclaimed or treated wastewater), and permits can cover noise and diesel backup generators. Approval could also require a
published catchment water balance, the water source and its distance, and a public hearing, tests that the Visakhapatnam record
lacks and that the NITI Aayog framework now proposes (Section~\ref{sec:siting}).
\textbf{Regionally}, grid planning and fair tariffs can stop households paying for data-centre load, and water can be
allocated by river basin, so that a city's supply and a power plant's cooling are judged together.
\textbf{Globally}, mandatory disclosure of energy, water and location-based carbon (from the actual local grid, not
offset by certificates bought elsewhere), site by site, would close most of the data gaps listed above, and
supply-chain standards can address mineral concentration and the e-waste collection gap.

\subsection{Conclusion}
The cloud era showed that computing can grow without much more electricity, as long as buildings keep getting more
efficient. That margin ran out around 2014, and the AI build-out arrived after it: compute took off in 2011, electricity
followed after 2018. Its footprint is still small against global totals (1.5\% of electricity, under 0.1\% of the
freshwater boundary), and that is exactly why global totals are the wrong yardstick. Weighted by size, a third of AI
capacity already sits in water-stressed basins and nearly half of the 2030 plan could; most of its water is evaporated
at power plants; and the capacity itself is concentrated in a handful of countries, with a Gini of 0.94. AI does not
have one footprint. It has a local, a regional and a global one, and each needs its own response.
